\section{Belief states and information-gathering actions}
\label{sec:C10:beliefs}
Under partial observability, a belief $b_t(e)=P(E_t=e\mid H_t)$ summarizes uncertainty in an ideal known model. After action $a$ and observation $o$, Bayes filtering gives $b_{t+1}(e')\propto P(o\mid e',a)\sum_eP(e'\mid e,a)b_t(e)$. Learned agents often use an approximate textual or latent state instead of an explicit distribution. That approximation should preserve which facts are observed, inferred, disputed, or unknown. A summary that converts uncertainty into an unqualified assertion changes the planning problem rather than merely compressing it.

Some actions primarily reduce uncertainty. Reading a file, querying transaction status, inspecting a browser element, or running a diagnostic can improve later decisions without immediately advancing the visible goal. Their value depends on what choices they can change. If two possible states require the same next action, distinguishing them has little immediate decision value; if they require opposite irreversible operations, the information can be decisive. This is the operational connection between representation sufficiency in Chapter~\ref{ch:C05} and planning. The planner needs distinctions relevant to its action set, not an exhaustive description of every environmental detail.

Information-gathering can itself have side effects. A diagnostic test can modify files or consume a rate-limited resource. A read can expose sensitive information to a model or external service. A browser refresh can discard an unsaved form. The planner should therefore model observation actions with costs and permissions rather than treating them as free oracle access. A common evaluation error gives a research agent a privileged state-inspection tool unavailable to the deployed system, then attributes the resulting gain entirely to better reasoning.

The belief update should incorporate negative evidence carefully. A missing response does not necessarily mean the operation failed; a missing memory result does not mean the underlying fact is false. The planner should distinguish absence of evidence from evidence of absence using the source's completeness contract. An authoritative database lookup can establish that no matching record exists under a defined query, while an approximate retrieval miss cannot. This distinction directly affects whether the next action should retry, search differently, or proceed with a state-changing operation.

